%% annex/p08/substructural.tex --- Part 8.6 proof material
%% Inherited from v1.1-verified lines 14001--14014.  The source gives a
%% sketch only; no proof-length ratio applies and no exponential rule is used.

\begin{annexitem}{annex:p8.imll.1}{entry:p08-imll}{IMLL soundness, completeness, and cut elimination}
\sketchnote{This is the inherited proof sketch for the multiplicative
intuitionistic core only.  It omits every rule-by-rule soundness case; the
well-definedness of the Lindenbaum quotient; verification of its induced
order, monoid laws, monotonicity, both directions of residuation, and truth
lemma; and the complete cut-elimination induction, including every principal
reduction, commuting conversion, and context split.  None of the four bang
rules on the separate exponential-extension sheet is covered.}

For soundness, interpret a finite antecedent multiset by the tensor of its
values, using \(1\) for the empty tensor.  Identity is reflexivity of the
order, cut is transitivity together with monotonicity of tensor, the tensor
and unit rules use the commutative-monoid laws, and the two linear-implication
rules use the displayed residuation equivalence.  Checking those assertions
for every rule and every context decomposition is omitted.

For completeness, order formulas by derivability and identify formulas that
derive one another.  Tensor, the tensor unit, and linear implication are to
induce operations on the resulting quotient.  The omitted Lindenbaum
argument must prove that these operations are well defined, that the induced
relation is a partial order, that tensor is a monotone commutative monoid
operation, and that both implications in

\[
 \lbrack A\rbrack\otimes\lbrack B\rbrack
 \leq\lbrack C\rbrack
 \quad\text{ iff }\quad
 \lbrack B\rbrack\leq\lbrack A\multimap C\rbrack
\]

hold.  The omitted truth lemma then identifies the value of each formula with
its quotient class.  A nonderivable sequent fails in this canonical algebra,
which gives the advertised completeness reduction once those steps have been
supplied.

For cut elimination, induct on the cut formula and on the heights of the two
premiss derivations.  A principal cut is reduced using the introduction rules
for its outermost connective; a nonprincipal cut is commuted above the last
inference on the appropriate side.  The omitted proof must enumerate the
tensor, linear-implication, and unit principal cases, all commuting cases, and
all multiset context splittings, and must verify a strict decrease in the
chosen induction measure.  The source provides this strategy but not those
cases, so the result remains sketched rather than proved.
\end{annexitem}
